% chapters/bab2-tinjauan-pustaka.tex -- BAB 2 TINJAUAN PUSTAKA
% Contoh teks ilmiah di bawah ini adalah placeholder untuk memperlihatkan
% hasil format (thesis.md); ganti dengan isi penelitian yang sebenarnya.

\chapter{TINJAUAN PUSTAKA}

\section{Penelitian Terdahulu}
Beberapa penelitian sebelumnya telah membahas klasifikasi sentimen
pada teks berbahasa Indonesia dengan berbagai pendekatan. Penelitian
oleh \citet{example2021} menerapkan metode Support Vector Machine
untuk mengklasifikasikan sentimen ulasan film, dan memperoleh akurasi
yang cukup tinggi meskipun masih terdapat kesalahan klasifikasi pada
kalimat yang mengandung negasi. Penelitian lain oleh
\citet{example2022} membandingkan beberapa arsitektur jaringan saraf
tiruan untuk klasifikasi sentimen ulasan produk, dan menyimpulkan
bahwa model berbasis RNN secara konsisten mengungguli model berbasis
Bag-of-Words pada dataset berukuran besar.

Penelitian oleh \citet{example2023} secara khusus mengkaji pengaruh
tahapan praproses, seperti normalisasi kata tidak baku dan penghapusan
stopword, terhadap kinerja model klasifikasi sentimen berbahasa
Indonesia. Hasil penelitian tersebut menunjukkan bahwa normalisasi
kata tidak baku memberikan kontribusi peningkatan akurasi yang lebih
besar dibandingkan dengan penghapusan stopword semata. Temuan ini
menjadi salah satu pertimbangan dalam merancang tahapan praproses
pada penelitian ini.

Selain isu praproses, kualitas dan cakupan data uji turut memengaruhi
keandalan model. Sebuah tolok ukur (benchmark) berskala besar untuk
klasifikasi teks pernah dievaluasi oleh \citet{example2016techreport},
yang menemukan bahwa performa model sangat bergantung pada
representasi fitur yang digunakan \citep{example2018incollection}.
Temuan tersebut sejalan dengan hasil sebelumnya yang menunjukkan
pentingnya tahapan praproses yang tepat sebelum data digunakan untuk
melatih model \citep{example2016techreport}.

Berbeda dengan penelitian-penelitian tersebut, penelitian ini berfokus
pada penerapan arsitektur LSTM secara spesifik untuk ulasan produk
berbahasa Indonesia dengan skema praproses yang disesuaikan, serta
melakukan perbandingan kinerja secara langsung dengan metode
pembelajaran mesin konvensional pada dataset yang sama
\citep{example2021,soemanto1993}.

\section{Landasan Teori}

\subsection{Pemrosesan bahasa alami dan analisis sentimen}
Pemrosesan Bahasa Alami (Natural Language Processing, NLP) merupakan
cabang ilmu komputer yang mempelajari interaksi antara komputer dan
bahasa manusia, baik dalam bentuk teks maupun ucapan. Salah satu
tugas populer dalam NLP adalah analisis sentimen (sentiment
analysis), yaitu proses identifikasi dan ekstraksi opini subjektif
dari suatu teks untuk menentukan polaritasnya, apakah bernilai
positif, negatif, atau netral. Analisis sentimen umumnya melibatkan
beberapa tahapan, yaitu pengumpulan data, praproses teks, ekstraksi
fitur, dan klasifikasi.

\subsection{Praproses teks berbahasa Indonesia}
Tahapan praproses teks bertujuan untuk mengubah teks mentah menjadi
representasi yang lebih terstruktur dan siap diolah oleh model
pembelajaran mesin. Tahapan yang umum dilakukan meliputi \textit{case
folding}, tokenisasi, normalisasi kata tidak baku, penghapusan
stopword, dan \textit{stemming}. Untuk bahasa Indonesia, tahapan
normalisasi kata tidak baku memegang peranan penting karena banyaknya
variasi ejaan informal yang digunakan pengguna pada media sosial dan
platform ulasan produk, misalnya kata ``bgs'' untuk ``bagus'' atau
``gak'' untuk ``tidak''.

\subsection{Long short-term memory}
Long Short-Term Memory (LSTM) adalah salah satu varian dari Recurrent
Neural Network (RNN) yang dirancang untuk mengatasi masalah
\textit{vanishing gradient} yang umum terjadi pada RNN standar ketika
memproses data sekuensial yang panjang. LSTM memperkenalkan struktur
sel memori (\textit{cell state}) yang dilengkapi dengan tiga gerbang
utama, yaitu \textit{forget gate}, \textit{input gate}, dan
\textit{output gate}. Ketiga gerbang tersebut bekerja sama untuk
mengatur informasi apa yang perlu dipertahankan, ditambahkan, atau
dibuang dari sel memori pada setiap langkah waktu, sehingga model
dapat mempelajari dependensi jangka panjang pada data teks dengan
lebih baik dibandingkan RNN standar.

Persamaan umum untuk memperbarui sel memori pada LSTM dapat dituliskan
sebagai berikut.
\begin{equation}
  c_t = f_t \odot c_{t-1} + i_t \odot \tilde{c}_t
  \label{eq:lstm-cell}
\end{equation}
dengan $f_t$ adalah nilai \textit{forget gate}, $i_t$ adalah nilai
\textit{input gate}, $c_{t-1}$ adalah sel memori pada langkah waktu
sebelumnya, dan $\tilde{c}_t$ adalah kandidat nilai sel memori baru
pada langkah waktu $t$, seperti pada Persamaan~\eqref{eq:lstm-cell}.

\section{Kerangka Berpikir}
Berdasarkan tinjauan pustaka di atas, penelitian ini menyusun
kerangka berpikir yang terdiri atas empat tahap utama, yaitu
pengumpulan data ulasan produk, praproses teks yang disesuaikan
dengan karakteristik bahasa Indonesia, perancangan dan pelatihan
model LSTM, serta evaluasi kinerja model dengan membandingkannya
terhadap metode pembelajaran mesin konvensional. Kerangka berpikir
ini menjadi dasar penyusunan metodologi penelitian pada bab
berikutnya.
